For the creation of directories, \LaTeX\ provides the standard commands
according to table~\ref{tab:final:toclists}. That is all the author has to do.
The classification of the directories in the outline of the work is explained
in the section~\ref{sec:content:section}.

\begin{table}
\centering
\begin{tabular}{ll}
\verb+\tableofcentents+     & Table of contents \\
\verb+\listoffigures+       & List of figures \\
\verb+\listoftables+        & List of tables \\
\verb+\lstlistoflistings+   & List of all source code \\
\end{tabular}
\caption{Directory commands}
\label{tab:final:toclists}
\end{table}

For definitions, sentences or self-defined sliding environments there exist
further commands. Please refer to the documentation of the corresponding
packages.\todo[inline]{search and link packages}.

Normally, only numbered sections appear in the table of contents.
However, if an unnumbered item is to be included in the table of contents, it
can be done with the \verb+\addcontentsline+ command, as in Listing~\ref{lst:final:addcontentsline}.
The command must always be called after the corresponding outline command.
Otherwise the page references will not be correct.

\lstinputlisting[language={[LaTeX]{TeX}},style=block,escapechar=\!,float,caption={adding additional directory entries},label={lst:final:addcontentsline},morekeywords={chapter,subsection}]{code/final_addcontentsline.en.tex}
